\section*{Chapter \ref{ch:g7:households} --- Households: Income, Spending and Saving}

\begin{solution}{exo:g7:households:1}
$42.7\,\%$ in the lowest fifth and $30.3\,\%$ in the highest.
\end{solution}

\begin{solution}{exo:g7:households:2}
$24\,000 - 2\,500 + 1\,800 = 23\,300$ dalens.
\end{solution}

\begin{solution}{exo:g7:households:3}
Food $5\,000 \div 20\,000 = 25\,\%$; housing $8\,000 \div 20\,000 = 40\,\%$; transport
$3\,000 \div 20\,000 = 15\,\%$; other $4\,000 \div 20\,000 = 20\,\%$. They add up to $100\,\%$.
\end{solution}

\begin{solution}{exo:g7:households:4}
$31\,000 - 28\,500 = 2\,500$ dalens.
\end{solution}

\begin{solution}{exo:g7:households:5}
A: $4\,800 \div 16\,000 = 30\,\%$. B: $5\,000 \div 25\,000 = 20\,\%$. C: $6\,000 \div 50\,000 =
12\,\%$. Yes: food spending rises with \omterm{def:g7:households:income}{income} ($4\,800$, $5\,000$, $6\,000$) while the \omterm{def:g6:food-and-energy:food-share}{food share}
falls.
\end{solution}

\begin{solution}{exo:g7:households:6}
$27.4 - 13.2 = 14.2$ percentage points. Food spending rises by $5\,250 - 3\,990 = 1\,260$ dalens.
\end{solution}

\begin{solution}{exo:g7:households:7}
$25.6 \div 14.6 \approx 1.75$: less than twice as large. Half of $27.4\,\%$ is $13.7\,\%$; the
\omterm{def:g6:food-and-energy:food-share}{food share} of $18.8\,\%$ is more than half as large.
\end{solution}

\begin{solution}{exo:g7:households:8}
It saves $2\,400 - 2\,250 = 150$ dalens a month, so $150 \times 12 = 1\,800$ dalens a year. The
part saved is $150 \div 2\,400 = 0.0625$, that is $6.25\,\%$.
\end{solution}

\begin{solution}{exo:g7:households:9}
For example: on average, \omterm{def:g7:households:household}{households} with higher \omterm{def:g7:households:income}{incomes} spending a larger share of their \omterm{def:g7:households:income}{income} on
food than \omterm{def:g7:households:household}{households} with lower \omterm{def:g7:households:income}{incomes}; or richer \omterm{def:g7:households:household}{households} spending less on food in \omterm{def:g7:money:money}{money} than
poorer ones.
\end{solution}

\begin{solution}{exo:g7:households:10}
(a) $5\,400 \div 30\,000 = 18\,\%$; $8\,500 \div 30\,000 \approx 0.283$, that is $28.3\,\%$.
(b) $5\,400 \div 2 = 2\,700$ dalens and $8\,500 \div 5 = 1\,700$ dalens per person. (c) The same
\omterm{def:g7:households:income}{income} feeds more people in a large \omterm{def:g7:households:household}{household}: per person it is poorer. A comparison of \omterm{def:g6:food-and-energy:food-share}{food
shares} across \omterm{def:g7:households:household}{households} of different sizes mixes the effect of \omterm{def:g7:households:income}{income} with the effect of size.
\end{solution}

\begin{solution}{exo:g7:households:11}
The first spends $0.30 \times 1\,000 = 300$ dalens more on food, the second $0.05 \times 1\,000 =
50$ dalens. Both spend more on food when their \omterm{def:g7:households:income}{income} rises; the richer \omterm{def:g7:households:household}{household} puts a smaller
part of the extra \omterm{def:g7:money:money}{money} into food, which is why the \omterm{def:g6:food-and-energy:food-share}{food share} falls as \omterm{def:g7:households:income}{income} rises.
\end{solution}

\begin{solution}{exo:g7:households:12}
The chart compares different \omterm{def:g7:households:household}{households}: the richest fifth may differ from the poorest in age, size,
place and tastes, not only in \omterm{def:g7:households:income}{income}. To see how one \omterm{def:g7:households:household}{household}'s \omterm{def:g7:households:budget}{budget} changes with its own
\omterm{def:g7:households:income}{income}, one needs the same \omterm{def:g7:households:household}{households} surveyed in several years.
\end{solution}

\begin{solution}{pb:g7:households:1}
\textbf{1.} A $25\,020 - 2\,600 + 2\,400 = 24\,820$; B $17\,310 - 1\,060 = 16\,250$; C $19\,400 -
1\,480 + 3\,600 = 21\,520$; D $36\,290 - 4\,860 + 1\,200 = 32\,630$; E $77\,600 - 18\,760 = 58\,840$
dalens.
\textbf{2.} C receives the largest \omterm{def:g7:households:transfer}{transfers} ($3\,600$); E pays the largest taxes ($18\,760$).
\textbf{3.} $18\,760 \div 77\,600 \approx 0.242$, that is $24.2\,\%$.
\textbf{4.} B ($16\,250$), C ($21\,520$), A ($24\,820$), D ($32\,630$), E ($58\,840$).
\textbf{5.} A $24\,820 - 24\,080 = 740$; B $0$; C $21\,520 - 21\,200 = 320$; D $32\,630 - 29\,440 =
3\,190$; E $58\,840 - 53\,060 = 5\,780$ dalens.
\textbf{6.} A $5\,490 \div 24\,820 \approx 22.1\,\%$; B $4\,160 \div 16\,250 = 25.6\,\%$; C $6\,890
\div 21\,520 \approx 32.0\,\%$; D $5\,390 \div 32\,630 \approx 16.5\,\%$; E $4\,080 \div 58\,840
\approx 6.9\,\%$.
\textbf{7.} B.
\textbf{8.} D: $3\,190 \div 32\,630 \approx 9.8\,\%$. E: $5\,780 \div 58\,840 \approx 9.8\,\%$.
\textbf{9.} Mostly: in the order B, C, A, D, E of \omterm{def:g7:households:disposable-income}{disposable income}, the \omterm{def:g6:food-and-energy:food-share}{food shares} are $25.6$,
$32.0$, $22.1$, $16.5$ and $6.9\,\%$. C departs from it: it has more \omterm{def:g7:households:income}{income} than B but a higher \omterm{def:g6:food-and-energy:food-share}{food
share}. It has five members, the most of the five \omterm{def:g7:households:household}{households}.
\textbf{10.} C: $6\,890 \div 5 = 1\,378$ dalens; E: $4\,080 \div 1 = 4\,080$ dalens.
\textbf{11.} \omterm{def:g7:households:income}{Income}: $39.6 \div 14.6 \approx 2.7$ times. Food spending: $5\,250 \div 3\,990 \approx
1.3$ times.
\textbf{12.} From the bottom quarter to the top, mean \omterm{def:g7:households:disposable-income}{disposable income} is about $2.7$ times
larger; food spending rises, from $3\,990$ to $5\,250$ dalens, but only about $1.3$ times; so the
\omterm{def:g6:food-and-energy:food-share}{food share} falls, from $27.4\,\%$ to $13.2\,\%$.
\end{solution}
